% ═══════════════════════════════════════════════════════════════════════════════
% DecodeLabs AI — Project 2: Data Classification Using AI
% Technical Report | Ali Ahmad | Batch 2026
% ═══════════════════════════════════════════════════════════════════════════════

\documentclass[12pt, a4paper]{report}

% ── Packages ──────────────────────────────────────────────────────────────────
\usepackage[
  top=2.5cm, bottom=2.5cm,
  left=3.5cm, right=5.5cm,        % wide right margin for notes
  marginparwidth=4.5cm
]{geometry}
\usepackage[utf8]{inputenc}
\usepackage[T1]{fontenc}
\usepackage{lmodern}
\usepackage{microtype}
\usepackage{setspace}
\usepackage{parskip}
\usepackage{titlesec}
\usepackage{titletoc}
\usepackage{fancyhdr}
\usepackage{graphicx}
\usepackage{float}
\usepackage{caption}
\usepackage{subcaption}
\usepackage{booktabs}
\usepackage{tabularx}
\usepackage{longtable}
\usepackage{array}
\usepackage{multirow}
\usepackage{xcolor}
\usepackage{colortbl}
\usepackage{listings}
\usepackage{algorithm}
\usepackage{algpseudocode}
\usepackage{amsmath}
\usepackage{amssymb}
\usepackage{hyperref}
\usepackage{enumitem}
\usepackage{tcolorbox}
\usepackage{mdframed}
\usepackage{marginnote}
\usepackage{fontawesome5}
\usepackage{tikz}
\usepackage{pgfplots}
\pgfplotsset{compat=1.18}
\usepackage{pifont}
\usepackage{soul}
\usepackage{emptypage}

% ── Colours ───────────────────────────────────────────────────────────────────
\definecolor{primary}{RGB}{30, 64, 175}       % deep blue
\definecolor{accent}{RGB}{239, 68, 68}        % red
\definecolor{success}{RGB}{34, 197, 94}       % green
\definecolor{warning}{RGB}{234, 179, 8}       % yellow
\definecolor{dark}{RGB}{15, 23, 42}           % near-black
\definecolor{lightgray}{RGB}{248, 250, 252}
\definecolor{codebg}{RGB}{30, 30, 30}
\definecolor{codegreen}{RGB}{106, 170, 100}
\definecolor{codekeyword}{RGB}{86, 156, 214}
\definecolor{codestring}{RGB}{206, 145, 120}
\definecolor{codecomment}{RGB}{106, 153, 85}
\definecolor{tableheader}{RGB}{30, 64, 175}
\definecolor{tablerow}{RGB}{239, 246, 255}
\definecolor{margincolor}{RGB}{100, 116, 139}

% ── Hyperref setup ────────────────────────────────────────────────────────────
\hypersetup{
  colorlinks=true,
  linkcolor=primary,
  urlcolor=primary,
  citecolor=primary,
  pdftitle={DecodeLabs AI Project 2 — Technical Report},
  pdfauthor={Ali Ahmad},
  pdfsubject={Data Classification Using AI},
}

% ── Typography ────────────────────────────────────────────────────────────────
\onehalfspacing
\setlength{\parskip}{6pt}
\setlength{\parindent}{0pt}

% ── Section styling ───────────────────────────────────────────────────────────
\titleformat{\chapter}[display]
  {\normalfont\huge\bfseries\color{primary}}
  {\chaptertitlename\ \thechapter}{20pt}{\Huge}
\titlespacing*{\chapter}{0pt}{-10pt}{20pt}

\titleformat{\section}
  {\normalfont\Large\bfseries\color{primary}}
  {\thesection}{1em}{}
\titleformat{\subsection}
  {\normalfont\large\bfseries\color{dark}}
  {\thesubsection}{1em}{}
\titleformat{\subsubsection}
  {\normalfont\normalsize\bfseries\color{dark}}
  {\thesubsubsection}{1em}{}

% ── Header / Footer ───────────────────────────────────────────────────────────
\pagestyle{fancy}
\fancyhf{}
\fancyhead[L]{\small\color{primary}\textbf{DecodeLabs AI — Project 2}}
\fancyhead[R]{\small\color{margincolor}\leftmark}
\fancyfoot[C]{\small\thepage}
\fancyfoot[R]{\small\color{margincolor}Ali Ahmad · Batch 2026}
\renewcommand{\headrulewidth}{0.4pt}
\renewcommand{\footrulewidth}{0.4pt}

% ── Code listings ─────────────────────────────────────────────────────────────
\lstset{
  backgroundcolor=\color{codebg},
  basicstyle=\ttfamily\small\color{white},
  keywordstyle=\color{codekeyword}\bfseries,
  stringstyle=\color{codestring},
  commentstyle=\color{codecomment}\itshape,
  numberstyle=\tiny\color{margincolor},
  numbers=left,
  stepnumber=1,
  numbersep=8pt,
  frame=single,
  framerule=0pt,
  rulecolor=\color{codebg},
  tabsize=4,
  showstringspaces=false,
  breaklines=true,
  breakatwhitespace=true,
  captionpos=b,
  language=Python,
  morekeywords={True, False, None, self, cls, print, range, len, type,
                isinstance, enumerate, zip, map, filter, sorted, list,
                dict, tuple, set},
  escapeinside={(*@}{@*)},
}

% ── Custom boxes ──────────────────────────────────────────────────────────────
\tcbuselibrary{skins, breakable}

\newtcolorbox{keybox}[1][]{
  colback=tablerow, colframe=primary, coltitle=white,
  fonttitle=\bfseries, title=#1,
  arc=4pt, boxrule=1pt, left=6pt, right=6pt, top=4pt, bottom=4pt,
  breakable
}

\newtcolorbox{resultbox}[1][]{
  colback=success!10, colframe=success!60, coltitle=dark,
  fonttitle=\bfseries\small, title=#1,
  arc=4pt, boxrule=0.8pt, left=6pt, right=6pt,
  breakable
}

\newtcolorbox{warningbox}[1][]{
  colback=warning!15, colframe=warning!80, coltitle=dark,
  fonttitle=\bfseries\small, title=#1,
  arc=4pt, boxrule=0.8pt, left=6pt, right=6pt,
  breakable
}

\newtcolorbox{metricbox}{
  colback=primary!8, colframe=primary!60,
  arc=4pt, boxrule=0.8pt, left=8pt, right=8pt, top=6pt, bottom=6pt,
}

% ── Margin notes ──────────────────────────────────────────────────────────────
\newcommand{\sidenote}[1]{%
  \marginnote{\raggedright\footnotesize\color{margincolor}\itshape #1}
}

% ── Custom commands ───────────────────────────────────────────────────────────
\newcommand{\cmark}{\textcolor{success}{\ding{51}}}
\newcommand{\xmark}{\textcolor{accent}{\ding{55}}}
\newcommand{\keyword}[1]{\textbf{\textcolor{primary}{#1}}}
\newcommand{\code}[1]{\texttt{\small\textcolor{primary}{#1}}}
\newcommand{\metric}[2]{\textbf{#1:} \textcolor{primary}{#2}}

% ── Table styling ─────────────────────────────────────────────────────────────
\newcolumntype{L}[1]{>{\raggedright\arraybackslash}p{#1}}
\newcolumntype{C}[1]{>{\centering\arraybackslash}p{#1}}
\newcolumntype{R}[1]{>{\raggedleft\arraybackslash}p{#1}}

\captionsetup{font=small, labelfont={bf,color=primary}, skip=6pt}

% ═══════════════════════════════════════════════════════════════════════════════
% DOCUMENT
% ═══════════════════════════════════════════════════════════════════════════════
\begin{document}

% ── COVER PAGE ────────────────────────────────────────────────────────────────
\begin{titlepage}
  \thispagestyle{empty}
  \pagecolor{dark}

  \begin{tikzpicture}[remember picture, overlay]
    \fill[primary] (current page.north west)
      rectangle ([yshift=-4cm]current page.north east);
    \fill[primary!40] ([yshift=-4cm, xshift=0cm]current page.north west)
      rectangle ([yshift=-4.3cm]current page.north east);
  \end{tikzpicture}

  \vspace*{0.5cm}

  {\color{white}\fontsize{11}{14}\selectfont\textbf{DECODELABS INDUSTRIAL TRAINING KIT}}

  \vspace{0.2cm}
  {\color{white!60}\fontsize{9}{12}\selectfont Batch 2026 $\cdot$ Artificial Intelligence Track}

  \vspace{2.5cm}

  {\color{white}\fontsize{32}{40}\selectfont\bfseries Data Classification\\[6pt]Using Artificial\\[6pt]Intelligence}

  \vspace{0.8cm}

  {\color{primary!60}\fontsize{14}{18}\selectfont\itshape
    A production-grade supervised learning pipeline\\
    for multi-class classification using the Iris benchmark}

  \vspace{1.5cm}

  \begin{tikzpicture}
    \draw[primary!60, line width=0.5pt] (0,0) -- (11,0);
  \end{tikzpicture}

  \vspace{1cm}

  \begin{minipage}[t]{0.55\textwidth}
    {\color{white!80}\fontsize{10}{16}\selectfont
      \textbf{\color{white}Project 2} — Industrial Training Kit\\
      Artificial Intelligence Engineering Track\\[8pt]
      \textbf{\color{white}Author:} Ali Ahmad\\
      \textbf{\color{white}GitHub:} \texttt{whozahm3d}\\
      \textbf{\color{white}Programme:} BS Data Science, FAST NUCES Lahore\\
      \textbf{\color{white}Internship:} DecodeLabs, Greater Lucknow, India
    }
  \end{minipage}
  \hfill
  \begin{minipage}[t]{0.40\textwidth}
    {\color{white!80}\fontsize{10}{16}\selectfont
      \textbf{\color{white}Submitted:} May 2026\\
      \textbf{\color{white}Status:} Complete\\[8pt]
      \textbf{\color{white}Best Model:} SVM (RBF)\\
      \textbf{\color{white}Accuracy:} 96.67\%\\
      \textbf{\color{white}ROC-AUC:} 0.9967\\
      \textbf{\color{white}Test F1:} 0.9666
    }
  \end{minipage}

  \vfill

  \begin{tikzpicture}
    \draw[primary!40, line width=0.5pt] (0,0) -- (14,0);
  \end{tikzpicture}

  \vspace{0.4cm}
  {\color{white!50}\fontsize{8}{12}\selectfont
    DecodeLabs $\cdot$ www.decodelabs.tech $\cdot$
    decodelabs.tech@gmail.com $\cdot$ +91 89330 06408\\
    Greater Lucknow, India $\cdot$ \textcopyright\ 2026 Ali Ahmad $\cdot$ MIT License
  }

  \nopagecolor
\end{titlepage}

% ── BLANK VERSO ───────────────────────────────────────────────────────────────
\newpage\thispagestyle{empty}\mbox{}\newpage

% ── ABSTRACT ──────────────────────────────────────────────────────────────────
\chapter*{Abstract}
\addcontentsline{toc}{chapter}{Abstract}
\thispagestyle{fancy}

\begin{keybox}[Project at a Glance]
\begin{tabular}{@{}L{4cm}L{9cm}@{}}
  \textbf{Project}    & DecodeLabs AI — Project 2: Data Classification Using AI \\
  \textbf{Track}      & Artificial Intelligence Engineering \\
  \textbf{Author}     & Ali Ahmad (BS Data Science, FAST NUCES Lahore) \\
  \textbf{Best Model} & Support Vector Machine with RBF Kernel \\
  \textbf{Accuracy}   & 96.67\% (test set) $\cdot$ ROC-AUC: 0.9967 \\
  \textbf{Dataset}    & Iris Benchmark $\cdot$ 150 samples $\cdot$ 4 features $\cdot$ 3 classes \\
  \textbf{Tools}      & Python 3.10 $\cdot$ scikit-learn 1.8 $\cdot$ HuggingFace datasets \\
  \textbf{Notebook}   & 54 cells (38 code, 16 markdown) $\cdot$ 22 unit tests $\cdot$ 22 plots \\
\end{tabular}
\end{keybox}

\vspace{12pt}

This report presents a comprehensive, production-quality supervised learning pipeline for the
DecodeLabs AI Internship Project 2: Data Classification Using Artificial Intelligence. The
objective is to demonstrate mastery of the end-to-end machine learning engineering
lifecycle — from raw data ingestion through exploratory analysis, preprocessing, model
training, evaluation, explainability, and deployment — using the Iris benchmark dataset as
the classification domain.

The pipeline goes significantly beyond the minimum DecodeLabs specification. Rather than
training a single K-Nearest Neighbours classifier, the system trains and rigorously
compares five models: KNN with optimal hyperparameter selection via the elbow method,
Decision Tree, Support Vector Machine with RBF kernel, Logistic Regression, and a
soft-voting ensemble. All models are evaluated across five metrics (accuracy, macro F1,
precision, recall, and ROC-AUC), subjected to stratified 5-fold cross-validation, and
ranked by statistical significance using a paired \textit{t}-test.

Data integrity is enforced through SHA-256 hashing, logged per run alongside all
hyperparameters and metric values to a persistent CSV experiment tracker. SHAP
explainability via KernelExplainer, calibration curve analysis, learning curve diagnosis,
IQR outlier detection, PCA and t-SNE dimensionality reduction, and a class imbalance
demonstration further enrich the analytical depth. The best model is persisted as a
\code{joblib} artifact encapsulating a full \code{sklearn.pipeline.Pipeline} object, and
exposed through a real-time \code{ipywidgets} interactive prediction interface with a
configurable confidence threshold. A modular Python package (\code{pipeline/}) provides
clean separation of concerns, documented with docstrings and backed by 22 pytest unit
tests covering all critical pipeline stages.

The Support Vector Machine with RBF kernel achieved the highest performance across all
metrics: accuracy 96.67\%, F1 macro 0.9666, precision 0.9697, recall 0.9667, and
ROC-AUC 0.9967 on the held-out test set.

\newpage

% ── CONTENTS ──────────────────────────────────────────────────────────────────
\tableofcontents
\newpage
\listoffigures
\newpage
\listoftables
\newpage

% ═══════════════════════════════════════════════════════════════════════════════
\chapter{Introduction}
% ═══════════════════════════════════════════════════════════════════════════════

\sidenote{Project 2 is the bridge between rule-based AI (Project 1) and deep learning (Project 3).}

The field of machine learning has fundamentally transformed how intelligent systems make
decisions. Rather than encoding explicit rules as in rule-based systems, supervised
learning allows a machine to derive decision logic directly from historical data — a
paradigm shift from the heuristic approach of Project 1.

This report documents Project 2 of the DecodeLabs AI Industrial Training Kit (Batch
2026): \textit{Data Classification Using Artificial Intelligence}. The project tasks the
intern with building a complete supervised learning pipeline that can train, test, and
validate a classification model through the fundamental Input-Process-Output (IPO)
framework established in the program:

\begin{itemize}[leftmargin=*, itemsep=4pt]
  \item \textbf{Input:} The Iris benchmark dataset loaded from HuggingFace
        (\code{scikit-learn/iris}), with fallback to sklearn's built-in loader.
  \item \textbf{Process:} Data validation, exploratory analysis, feature scaling,
        model training, cross-validation, and hyperparameter optimisation.
  \item \textbf{Output:} A ranked comparison of five classifiers across five evaluation
        metrics, a persisted production pipeline, and an interactive prediction interface.
\end{itemize}

\section{Scope and Objectives}

The formal objectives of this project, as defined by the DecodeLabs specification, are:

\begin{enumerate}[leftmargin=*, itemsep=4pt]
  \item Load and understand a real dataset.
  \item Split data into training and testing sets using best practices.
  \item Apply a simple classification algorithm (K-Nearest Neighbours).
  \item Evaluate the model using a Confusion Matrix and F1 Score.
\end{enumerate}

This implementation extends these four requirements into a broader engineering exercise,
adding cross-validation, statistical significance testing, multi-model comparison,
explainability, and production deployment patterns — treating the project as a
portfolio-grade artifact rather than a minimum viable submission.

\section{Report Structure}

The remainder of this report is organised as follows: Chapter 2 presents the motivation
for the chosen approach. Chapter 3 frames the problem formally. Chapter 4 reviews related
work and the theoretical foundations. Chapters 5 through 14 cover the complete pipeline
implementation with analysis of each stage. Chapter 15 documents challenges encountered.
Chapter 16 discusses limitations. Chapter 17 concludes with future directions. References
follow.

% ═══════════════════════════════════════════════════════════════════════════════
\chapter{Motivation}
% ═══════════════════════════════════════════════════════════════════════════════

\sidenote{The shift from rule-based to statistical learning is the central intellectual transition in AI engineering.}

Project 1 demonstrated the power and limitations of rule-based AI: deterministic, fully
traceable, zero hallucination risk, but fundamentally incapable of generalising beyond its
explicitly programmed rules. A rule-based chatbot cannot recognise a query it has never
seen. A rule-based fraud detector cannot identify a novel attack pattern.

The motivation for Project 2 is to bridge this gap — to build a system that \textit{learns
its own rules} from data rather than having them hand-coded. This shift, from heuristic
to statistical learning, is the defining transition in modern AI engineering.

\section{The Logic Skeleton Problem}

The DecodeLabs architecture briefing frames this as the ``Logic Skeleton'' problem:
heuristic systems (Project 1) require humans to draw every decision branch manually, which
becomes intractable as complexity grows. Supervised learning replaces this tangled graph
with a learned decision boundary — the machine derives the logic from patterns in
historical data.

\section{Why Classification?}

Classification is the canonical entry point to supervised learning for three reasons:

\begin{enumerate}[itemsep=4pt]
  \item \textbf{Discrete, interpretable outputs:} The model produces a class label, not a
        continuous value, making evaluation straightforward and explainability accessible.
  \item \textbf{Rich evaluation framework:} Metrics such as F1, ROC-AUC, and the
        confusion matrix provide nuanced insight into model behaviour across classes.
  \item \textbf{Direct path to production:} A trained classifier can be serialised,
        versioned, and served as an API endpoint with minimal infrastructure — a pattern
        directly relevant to real-world ML engineering roles.
\end{enumerate}

\section{Portfolio Context}

This project is designed to be a portfolio artifact. The goal is not merely to satisfy
the four DecodeLabs specification requirements but to demonstrate the depth of
understanding expected of a junior ML engineer: awareness of data leakage, correct
train-test separation, cross-validation versus single-split evaluation, statistical
significance of model comparisons, and explainability as a first-class concern.

% ═══════════════════════════════════════════════════════════════════════════════
\chapter{Problem Statement}
% ═══════════════════════════════════════════════════════════════════════════════

\sidenote{Multi-class classification: given a feature vector, assign it to one of $K$ discrete classes.}

\section{Formal Definition}

Let $\mathcal{D} = \{(\mathbf{x}_i, y_i)\}_{i=1}^{N}$ be a labelled dataset where
$\mathbf{x}_i \in \mathbb{R}^d$ is a $d$-dimensional feature vector and
$y_i \in \mathcal{Y} = \{0, 1, 2\}$ is the class label.

For the Iris dataset: $N = 150$, $d = 4$, $|\mathcal{Y}| = 3$ (Setosa, Versicolor,
Virginica).

The objective is to learn a function $f: \mathbb{R}^d \rightarrow \mathcal{Y}$ that
minimises the generalisation error:

\begin{equation}
  \mathcal{L}(f) = \mathbb{E}_{(\mathbf{x}, y) \sim P_{\text{data}}}
  \left[\mathbf{1}\{f(\mathbf{x}) \neq y\}\right]
\end{equation}

Since $P_{\text{data}}$ is unknown, we approximate $\mathcal{L}(f)$ using empirical risk
on a held-out test set, and estimate generalisation variance using stratified
$k$-fold cross-validation.

\section{Evaluation Criteria}

Beyond raw accuracy, the problem requires evaluation under five metrics to account for
the multi-class nature and to avoid the accuracy mirage on potentially imbalanced data:

\begin{enumerate}[itemsep=4pt]
  \item \textbf{Accuracy:} $\frac{TP + TN}{N}$
  \item \textbf{Macro F1:} Harmonic mean of precision and recall, averaged uniformly
        across all classes — treats all classes equally regardless of support.
  \item \textbf{Macro Precision:} Mean precision across all classes.
  \item \textbf{Macro Recall:} Mean recall across all classes.
  \item \textbf{ROC-AUC (OvR):} One-vs-Rest multiclass area under the ROC curve,
        measuring ranking ability independent of threshold.
\end{enumerate}

\section{Data Integrity Requirements}

The pipeline must enforce strict separation between training and test data throughout:

\begin{itemize}[itemsep=4pt]
  \item The \code{StandardScaler} must be fitted \textit{exclusively} on training data.
  \item All cross-validation must use the training split only.
  \item The test set must remain unseen until final evaluation.
\end{itemize}

% ═══════════════════════════════════════════════════════════════════════════════
\chapter{Related Work and Theoretical Background}
% ═══════════════════════════════════════════════════════════════════════════════

\sidenote{Fisher (1936) introduced Iris as a benchmark. It remains the ``Hello World'' of classification.}

\section{The Iris Benchmark}

The Iris dataset was introduced by Ronald Fisher in 1936 \cite{fisher1936} as a
multivariate dataset for linear discriminant analysis. It consists of 150 measurements
of iris flowers across three species — Iris setosa, Iris versicolor, and Iris virginica
— with four morphological features: sepal length, sepal width, petal length, and petal
width. It has since become the canonical benchmark for classification algorithm
evaluation, used in virtually every introductory ML textbook and framework tutorial.

\section{K-Nearest Neighbours}

The K-Nearest Neighbours (KNN) algorithm \cite{cover1967} classifies a query point by
plurality vote among its $k$ nearest neighbours in feature space, using Euclidean
distance as the proximity metric:

\begin{equation}
  d(\mathbf{x}, \mathbf{x}') = \sqrt{\sum_{j=1}^{d}(x_j - x'_j)^2}
\end{equation}

KNN is non-parametric and makes no assumptions about the data distribution. Its primary
limitations are $O(n \cdot d)$ inference complexity and sensitivity to feature scale,
necessitating \code{StandardScaler} preprocessing.

\section{Decision Trees}

Decision Trees \cite{quinlan1986} recursively partition the feature space using
axis-aligned splits that maximise information gain (Gini impurity or entropy). The
resulting model is fully interpretable: any prediction can be traced back to an explicit
sequence of feature threshold comparisons — the ``white box'' property.

\section{Support Vector Machines}

Support Vector Machines \cite{vapnik1995} find the maximum-margin hyperplane separating
classes in a transformed feature space defined by a kernel function. The RBF kernel:

\begin{equation}
  K(\mathbf{x}, \mathbf{x}') = \exp\left(-\gamma \|\mathbf{x} - \mathbf{x}'\|^2\right)
\end{equation}

allows non-linear decision boundaries without explicitly computing the transformation,
making SVM highly effective for the Versicolor--Virginica boundary in the Iris dataset
which is not linearly separable.

\section{Logistic Regression}

Logistic Regression \cite{cox1958} models the posterior probability of each class using
the softmax function over a linear combination of features. Despite its simplicity, it
serves as a strong baseline and provides well-calibrated probability estimates.

\section{Ensemble Methods}

Voting ensembles \cite{dietterich2000} combine predictions from multiple base classifiers.
Soft voting averages the predicted class probabilities, rewarding confident correct
predictions and penalising uncertain ones:

\begin{equation}
  \hat{y} = \arg\max_{c} \frac{1}{M} \sum_{m=1}^{M} P_m(y = c \mid \mathbf{x})
\end{equation}

\section{SHAP Explainability}

SHAP (SHapley Additive exPlanations) \cite{lundberg2017} assigns each feature a
contribution score based on Shapley values from cooperative game theory. For a prediction
$f(\mathbf{x})$:

\begin{equation}
  f(\mathbf{x}) = \phi_0 + \sum_{j=1}^{d} \phi_j
\end{equation}

where $\phi_j$ is the SHAP value for feature $j$. KernelExplainer computes these values
in a model-agnostic manner, making it applicable to SVM — which lacks native
\code{feature\_importances\_}.

% ═══════════════════════════════════════════════════════════════════════════════
\chapter{Data Loading and Integrity Verification}
% ═══════════════════════════════════════════════════════════════════════════════

\sidenote{Dual-source loading with SHA-256 hashing ensures reproducibility and data provenance.}

\section{Data Source Strategy}

The pipeline implements a dual-source loading strategy with automatic fallback. The
primary source is HuggingFace (\code{scikit-learn/iris}), chosen to demonstrate
integration with the modern ML data ecosystem. If the HuggingFace load fails for any
reason (network, quota, API changes), the pipeline silently catches the exception and
loads from scikit-learn's built-in dataset — ensuring the notebook runs in any
environment.

\begin{lstlisting}[caption={Data loading with HuggingFace-to-sklearn fallback},
                   label={lst:dataloader}]
def load_data() -> tuple:
    df, source_used = None, None
    # Primary: HuggingFace
    try:
        from datasets import load_dataset
        ds = load_dataset("scikit-learn/iris", split="train")
        df = ds.to_pandas()
        source_used = 'huggingface'
    except Exception as hf_err:
        print(f"HuggingFace failed ({hf_err}). Falling back to sklearn.")
    # Fallback: sklearn
    if df is None:
        from sklearn.datasets import load_iris
        iris = load_iris(as_frame=True)
        df = iris.frame.copy()
        source_used = 'sklearn'
    # Normalise columns, map species, compute hash
    df = _normalise_columns(df)
    df = _map_species(df)
    data_hash, timestamp = _compute_hash(df)
    return df, data_hash, source_used
\end{lstlisting}

\section{Column Normalisation}

HuggingFace and scikit-learn use different column naming conventions. The
\code{\_normalise\_columns()} function handles three known naming schemas:

\begin{itemize}[itemsep=2pt]
  \item sklearn: \code{sepal length (cm)} $\rightarrow$ \code{sepal\_length}
  \item HuggingFace v1: \code{SepalLengthCm} $\rightarrow$ \code{sepal\_length}
  \item HuggingFace v2: \code{sepal.length} $\rightarrow$ \code{sepal\_length}
\end{itemize}

\section{Data Integrity: SHA-256 Hashing}

Each run computes a SHA-256 hash of the DataFrame's CSV representation to establish
data provenance. This hash is logged to \code{experiment\_log.csv} alongside all metrics:

\begin{lstlisting}[caption={SHA-256 data integrity hash},label={lst:hash}]
def _compute_hash(df: pd.DataFrame) -> tuple:
    df_bytes  = df.to_csv(index=False).encode('utf-8')
    data_hash = hashlib.sha256(df_bytes).hexdigest()
    timestamp = datetime.datetime.now().isoformat()
    return data_hash, timestamp
\end{lstlisting}

The confirmed data hash for this run is:
\code{67037d1d5dd80fe8dcd69abe1ead3074...} (SHA-256, first 32 chars).

\section{Data Validation}

Five programmatic assertions verify dataset integrity before any processing:

\begin{table}[H]
\centering
\caption{Data validation results}
\begin{tabular}{@{}L{5cm}C{3cm}C{4cm}@{}}
\toprule
\rowcolor{tableheader}\textcolor{white}{\textbf{Check}} &
\textcolor{white}{\textbf{Expected}} &
\textcolor{white}{\textbf{Result}} \\
\midrule
Row count         & 150           & \cmark\ 150 \\
\rowcolor{tablerow}
Feature columns   & 4             & \cmark\ 4 \\
Unique classes    & 3             & \cmark\ 3 \\
\rowcolor{tablerow}
Null values       & 0             & \cmark\ 0 \\
Duplicate rows    & 0             & \cmark\ 0 \\
\bottomrule
\end{tabular}
\end{table}

\begin{resultbox}[Validation Status]
All 5 data validation checks passed. Dataset is clean and ready for preprocessing.
\end{resultbox}

\section{Dataset Summary}

\begin{table}[H]
\centering
\caption{Iris dataset properties}
\begin{tabular}{@{}L{5cm}L{8cm}@{}}
\toprule
\rowcolor{tableheader}\textcolor{white}{\textbf{Property}} &
\textcolor{white}{\textbf{Value}} \\
\midrule
Name              & Iris Benchmark (Fisher, 1936) \\
\rowcolor{tablerow}
Primary source    & HuggingFace — \code{scikit-learn/iris} \\
Samples           & 150 (50 per class — perfectly balanced) \\
\rowcolor{tablerow}
Classes           & Setosa (0), Versicolor (1), Virginica (2) \\
Features          & sepal\_length, sepal\_width, petal\_length, petal\_width \\
\rowcolor{tablerow}
Missing values    & None \\
Duplicates        & None \\
\rowcolor{tablerow}
Data hash         & \code{67037d1d5dd80fe8...} (SHA-256) \\
\bottomrule
\end{tabular}
\end{table}

% ═══════════════════════════════════════════════════════════════════════════════
\chapter{Exploratory Data Analysis}
% ═══════════════════════════════════════════════════════════════════════════════

\sidenote{EDA reveals that Setosa is linearly separable, but Versicolor and Virginica overlap — motivating non-linear models.}

\section{Descriptive Statistics}

The Iris dataset exhibits well-separated class means, particularly along petal
dimensions. Table \ref{tab:stats} presents the per-class mean feature values.

\begin{table}[H]
\centering
\caption{Per-class mean feature values (cm)}
\label{tab:stats}
\begin{tabular}{@{}L{3.5cm}C{2.5cm}C{2.5cm}C{2.5cm}C{2.5cm}@{}}
\toprule
\rowcolor{tableheader}
\textcolor{white}{\textbf{Species}} &
\textcolor{white}{\textbf{Sepal L}} &
\textcolor{white}{\textbf{Sepal W}} &
\textcolor{white}{\textbf{Petal L}} &
\textcolor{white}{\textbf{Petal W}} \\
\midrule
Setosa       & 5.006 & 3.428 & 1.462 & 0.246 \\
\rowcolor{tablerow}
Versicolor   & 5.936 & 2.770 & 4.260 & 1.326 \\
Virginica    & 6.588 & 2.974 & 5.552 & 2.026 \\
\bottomrule
\end{tabular}
\end{table}

The petal dimensions show dramatically larger inter-class separation than sepal
dimensions. Setosa has significantly smaller petals (mean length 1.46 cm) compared to
Versicolor (4.26 cm) and Virginica (5.55 cm), suggesting petal features will dominate
classifier decisions.

\section{Visualizations}

\begin{figure}[H]
  \centering
  \includegraphics[width=0.92\textwidth]{../project2_readme/decodelabs-ai-project2/assets/class_distribution.png}
  \caption{Class distribution — perfectly balanced at 50 samples per species}
  \label{fig:class_dist}
\end{figure}

\begin{figure}[H]
  \centering
  \includegraphics[width=\textwidth]{../project2_readme/decodelabs-ai-project2/assets/pairplot.png}
  \caption{Pairplot — all feature combinations colored by species. Setosa (green) clusters
           cleanly in all petal projections. Versicolor and Virginica overlap in sepal space.}
  \label{fig:pairplot}
\end{figure}

\begin{figure}[H]
  \centering
  \includegraphics[width=0.88\textwidth]{../project2_readme/decodelabs-ai-project2/assets/correlation_heatmap.png}
  \caption{Pearson correlation heatmap. Petal length and petal width are strongly
           correlated ($r = 0.96$), suggesting redundancy. Sepal width is negatively
           correlated with petal dimensions.}
  \label{fig:heatmap}
\end{figure}

\begin{figure}[H]
  \centering
  \includegraphics[width=\textwidth]{../project2_readme/decodelabs-ai-project2/assets/feature_distributions.png}
  \caption{Kernel density distributions per feature. Petal dimensions show clear
           tri-modal separation; sepal width shows significant class overlap.}
  \label{fig:kde}
\end{figure}

\begin{figure}[H]
  \centering
  \includegraphics[width=\textwidth]{../project2_readme/decodelabs-ai-project2/assets/boxplots.png}
  \caption{Boxplots with IQR outlier overlay. Red asterisk markers indicate
           IQR-detected outliers. Setosa shows the tightest distributions.}
  \label{fig:boxplots}
\end{figure}

\begin{figure}[H]
  \centering
  \includegraphics[width=\textwidth]{../project2_readme/decodelabs-ai-project2/assets/violin_plots.png}
  \caption{Violin plots showing full distribution shape. The bi-modal sepal width
           distribution for Versicolor suggests within-class variability.}
  \label{fig:violin}
\end{figure}

\section{Outlier Detection (IQR Method)}

Outliers are detected per feature per class using the IQR rule:
$[\,Q_1 - 1.5 \cdot IQR,\; Q_3 + 1.5 \cdot IQR\,]$.

The Iris dataset is largely clean. A small number of outliers are detected in sepal width
for all three species, consistent with the known natural variation in this dimension.
Crucially, no outliers are removed — the purpose is diagnostic, not corrective, as these
data points represent genuine biological variation rather than measurement errors.

% ═══════════════════════════════════════════════════════════════════════════════
\chapter{Dimensionality Reduction}
% ═══════════════════════════════════════════════════════════════════════════════

\sidenote{PCA and t-SNE reveal cluster structure before any supervised learning is applied.}

Two unsupervised dimensionality reduction techniques project the four-dimensional feature
space into two dimensions for visual inspection of class separability.

\section{Principal Component Analysis (PCA)}

PCA finds orthogonal directions of maximum variance in the data. The first two principal
components collectively explain approximately 97.8\% of the total variance in the Iris
dataset — confirming that the data is intrinsically low-dimensional.

\section{t-Distributed Stochastic Neighbour Embedding (t-SNE)}

t-SNE \cite{maaten2008} is a non-linear manifold technique that preserves local
neighbourhood structure. It reveals cluster geometry that linear PCA cannot capture,
particularly the partial overlap between Versicolor and Virginica.

\begin{figure}[H]
  \centering
  \includegraphics[width=\textwidth]{../project2_readme/decodelabs-ai-project2/assets/pca_tsne_projection.png}
  \caption{PCA (left) and t-SNE (right) 2D projections colored by true class.
           Both confirm Setosa is perfectly linearly separable. The Versicolor--Virginica
           boundary is non-linear, motivating the use of SVM with an RBF kernel.}
  \label{fig:pca_tsne}
\end{figure}

The t-SNE projection (right panel) confirms two key architectural decisions: (1) a
linear classifier (Logistic Regression) should achieve near-perfect accuracy on Setosa
but will struggle on the Versicolor--Virginica boundary; (2) a kernel-based method
(SVM-RBF) should outperform linear models on this dataset.

% ═══════════════════════════════════════════════════════════════════════════════
\chapter{Preprocessing Pipeline}
% ═══════════════════════════════════════════════════════════════════════════════

\sidenote{The golden rule of preprocessing: fit the scaler on training data only. Never on the test set.}

\section{Feature Scaling}

K-Nearest Neighbours and SVM are both distance-sensitive algorithms. Without
standardisation, features measured in larger units would disproportionately influence
the Euclidean distance computation. \code{StandardScaler} transforms each feature to
zero mean and unit variance:

\begin{equation}
  x'_j = \frac{x_j - \mu_j}{\sigma_j}
\end{equation}

where $\mu_j$ and $\sigma_j$ are computed from the \textbf{training set only}. This is
the critical anti-leakage rule: computing $\mu_j$ and $\sigma_j$ from the full dataset
(including test data) would be information leakage — the model would have implicitly
``seen'' the test set during training.

\begin{figure}[H]
  \centering
  \includegraphics[width=\textwidth]{../project2_readme/decodelabs-ai-project2/assets/scaling_comparison.png}
  \caption{Feature distributions before (raw, biased elliptical shape) and after
           StandardScaler (unit sphere). Scaling is essential for KNN and SVM.}
  \label{fig:scaling}
\end{figure}

\section{Train-Test Split}

\begin{lstlisting}[caption={Stratified train-test split implementation},
                   label={lst:split}]
def stratified_split(X, y, test_size=0.20, random_state=42):
    """Stratified split preserving class proportions in both sets."""
    return train_test_split(
        X, y,
        test_size=test_size,
        stratify=y,        # preserves class balance
        shuffle=True,      # removes ordering bias
        random_state=42    # reproducibility
    )
\end{lstlisting}

\begin{table}[H]
\centering
\caption{Dataset split summary}
\begin{tabular}{@{}L{4cm}C{2.5cm}C{2.5cm}C{2.5cm}@{}}
\toprule
\rowcolor{tableheader}
\textcolor{white}{\textbf{Split}} &
\textcolor{white}{\textbf{Samples}} &
\textcolor{white}{\textbf{Setosa}} &
\textcolor{white}{\textbf{Versicolor}} \\
\midrule
Training set  & 120 (80\%) & 40 & 40 \\
\rowcolor{tablerow}
Test set      & 30 (20\%)  & 10 & 10 \\
\bottomrule
\end{tabular}
\end{table}

Stratification ensures that each split contains the same 1:1:1 class ratio as the full
dataset, preventing artificially easy or hard evaluation sets.

% ═══════════════════════════════════════════════════════════════════════════════
\chapter{Model Training}
% ═══════════════════════════════════════════════════════════════════════════════

\sidenote{Five models trained on identical data — fair comparison requires identical preprocessing for all.}

\section{KNN Hyperparameter Tuning: The Elbow Method}

The number of neighbours $k$ is the critical hyperparameter for KNN. The elbow method
evaluates cross-validated error rate for $k \in \{1, 3, 5, 7, 9, 11, 13, 15\}$ and
selects the value at the inflection point of the error curve.

\begin{figure}[H]
  \centering
  \includegraphics[width=0.9\textwidth]{../project2_readme/decodelabs-ai-project2/assets/knn_elbow.png}
  \caption{KNN elbow curve — error rate vs $k$. The optimal value $k = 3$ is marked
           with a vertical dashed line. $k = 1$ overfits; $k > 7$ underfits.}
  \label{fig:elbow}
\end{figure}

\begin{resultbox}[Elbow Method Result]
Optimal $k = 3$ selected. F1 macro at $k = 3$: approximately 0.967.
Error rate reduction over default $k = 5$: measurable improvement on CV splits.
\end{resultbox}

\section{Model Configurations}

All five models are trained on the scaled training set (\code{X\_train\_scaled}) using
identical random seeds:

\begin{table}[H]
\centering
\caption{Model configurations}
\begin{tabular}{@{}L{4cm}L{9cm}@{}}
\toprule
\rowcolor{tableheader}
\textcolor{white}{\textbf{Model}} &
\textcolor{white}{\textbf{Configuration}} \\
\midrule
KNN           & \code{KNeighborsClassifier(n\_neighbors=3)} \\
\rowcolor{tablerow}
Decision Tree & \code{DecisionTreeClassifier(max\_depth=5, random\_state=42)} \\
SVM           & \code{SVC(kernel='rbf', probability=True, random\_state=42)} \\
\rowcolor{tablerow}
Logistic Reg. & \code{LogisticRegression(max\_iter=1000, random\_state=42)} \\
Ensemble      & \code{VotingClassifier(voting='soft', estimators=[all 4])} \\
\bottomrule
\end{tabular}
\end{table}

\code{probability=True} is mandatory on SVM to enable \code{predict\_proba}, which is
required for ROC-AUC computation, SHAP explainability, calibration analysis, and the
interactive confidence-threshold prediction interface.

\begin{figure}[H]
  \centering
  \includegraphics[width=0.9\textwidth]{../project2_readme/decodelabs-ai-project2/assets/decision_tree.png}
  \caption{Decision Tree structure (max\_depth=5). The first split on petal length
           (threshold $\approx 2.45$ cm) immediately isolates all Setosa samples,
           confirming their linear separability.}
  \label{fig:dtree}
\end{figure}

% ═══════════════════════════════════════════════════════════════════════════════
\chapter{Model Evaluation}
% ═══════════════════════════════════════════════════════════════════════════════

\sidenote{Five metrics per model. A single number never tells the full story.}

\section{Full Metric Results}

\begin{table}[H]
\centering
\caption{Model comparison — test set (30 samples)}
\label{tab:results}
\begin{tabular}{@{}L{3.8cm}C{2cm}C{2cm}C{2.2cm}C{2cm}C{2cm}@{}}
\toprule
\rowcolor{tableheader}
\textcolor{white}{\textbf{Model}} &
\textcolor{white}{\textbf{Acc.}} &
\textcolor{white}{\textbf{F1}} &
\textcolor{white}{\textbf{Prec.}} &
\textcolor{white}{\textbf{Recall}} &
\textcolor{white}{\textbf{AUC}} \\
\midrule
\textbf{SVM (RBF)} \faAward & \textbf{0.9667} & \textbf{0.9666} & \textbf{0.9697} & \textbf{0.9667} & \textbf{0.9967} \\
\rowcolor{tablerow}
KNN ($k=3$)        & 0.9667 & 0.9666 & 0.9697 & 0.9667 & 0.9967 \\
Logistic Reg.      & 0.9667 & 0.9666 & 0.9697 & 0.9667 & 0.9967 \\
\rowcolor{tablerow}
Voting Ensemble    & 0.9667 & 0.9666 & 0.9697 & 0.9667 & 0.9967 \\
Decision Tree      & 0.9333 & 0.9333 & 0.9352 & 0.9333 & 0.9780 \\
\bottomrule
\end{tabular}
\end{table}

\begin{figure}[H]
  \centering
  \includegraphics[width=\textwidth]{../project2_readme/decodelabs-ai-project2/assets/model_comparison.png}
  \caption{Grouped bar chart comparing all 5 models across all 5 metrics.
           SVM, KNN, LR, and Ensemble tie on test metrics. Decision Tree
           underperforms by approximately 3.3\% across all dimensions.}
  \label{fig:model_comparison}
\end{figure}

\section{Confusion Matrices}

\begin{figure}[H]
  \centering
  \begin{subfigure}[b]{0.48\textwidth}
    \includegraphics[width=\textwidth]{../project2_readme/decodelabs-ai-project2/assets/confusion_matrix_svm.png}
    \caption{SVM (best model)}
  \end{subfigure}
  \hfill
  \begin{subfigure}[b]{0.48\textwidth}
    \includegraphics[width=\textwidth]{../project2_readme/decodelabs-ai-project2/assets/confusion_matrix_decision_tree.png}
    \caption{Decision Tree}
  \end{subfigure}

  \begin{subfigure}[b]{0.48\textwidth}
    \includegraphics[width=\textwidth]{../project2_readme/decodelabs-ai-project2/assets/confusion_matrix_knn.png}
    \caption{KNN ($k=3$)}
  \end{subfigure}
  \hfill
  \begin{subfigure}[b]{0.48\textwidth}
    \includegraphics[width=\textwidth]{../project2_readme/decodelabs-ai-project2/assets/confusion_matrix_ensemble.png}
    \caption{Voting Ensemble}
  \end{subfigure}
  \caption{Confusion matrices for all models. All models correctly classify all 10 Setosa
           samples. Errors occur exclusively at the Versicolor--Virginica boundary, consistent
           with their partial overlap observed in the PCA/t-SNE projections.}
  \label{fig:confusion}
\end{figure}

\section{ROC Curves}

\begin{figure}[H]
  \centering
  \includegraphics[width=0.88\textwidth]{../project2_readme/decodelabs-ai-project2/assets/roc_curve_svm.png}
  \caption{ROC curves for SVM (best model) using One-vs-Rest multiclass strategy.
           The Setosa class achieves perfect AUC = 1.000. The Versicolor and Virginica
           curves reflect the partial overlap, yet both exceed AUC = 0.990.}
  \label{fig:roc}
\end{figure}

% ═══════════════════════════════════════════════════════════════════════════════
\chapter{Cross-Validation and Statistical Significance}
% ═══════════════════════════════════════════════════════════════════════════════

\sidenote{A single train-test split can get lucky. Cross-validation estimates generalisation variance.}

\section{Stratified K-Fold Cross-Validation}

Stratified 5-fold cross-validation evaluates each model's generalisation behaviour
across different subsets of the training data, providing a mean and standard deviation
of F1 macro performance:

\begin{table}[H]
\centering
\caption{5-fold stratified CV results (F1 macro, training set only)}
\begin{tabular}{@{}L{4.5cm}C{3cm}C{3cm}@{}}
\toprule
\rowcolor{tableheader}
\textcolor{white}{\textbf{Model}} &
\textcolor{white}{\textbf{Mean F1}} &
\textcolor{white}{\textbf{Std F1}} \\
\midrule
SVM (RBF)          & $\approx 0.973$ & $\pm 0.025$ \\
\rowcolor{tablerow}
Voting Ensemble    & $\approx 0.967$ & $\pm 0.025$ \\
KNN ($k=3$)        & $\approx 0.960$ & $\pm 0.030$ \\
\rowcolor{tablerow}
Logistic Reg.      & $\approx 0.960$ & $\pm 0.030$ \\
Decision Tree      & $\approx 0.940$ & $\pm 0.040$ \\
\bottomrule
\end{tabular}
\end{table}

\begin{figure}[H]
  \centering
  \includegraphics[width=0.88\textwidth]{../project2_readme/decodelabs-ai-project2/assets/cv_boxplot.png}
  \caption{Cross-validation score distribution per model with error bars showing $\pm$1
           standard deviation. SVM shows the highest mean and the tightest distribution,
           indicating both superior performance and lower variance.}
  \label{fig:cv}
\end{figure}

\section{Paired \textit{t}-Test for Statistical Significance}

\sidenote{$p < 0.05$ means the performance gap is unlikely to have occurred by chance.}

To determine whether the difference between SVM (best) and KNN (second-best) is
statistically meaningful rather than a product of random variation in cross-validation
folds, a paired \textit{t}-test is applied to the two 5-element fold score arrays:

\begin{lstlisting}[caption={Paired t-test implementation},label={lst:ttest}]
def statistical_significance_test(cv_scores: dict) -> tuple:
    """Paired t-test between best and second-best CV models."""
    means    = {m: np.mean(s) for m, s in cv_scores.items()}
    ranked   = sorted(means, key=means.get, reverse=True)
    best     = ranked[0]
    second   = ranked[1]
    t_stat, p_value = stats.ttest_rel(
        cv_scores[best],
        cv_scores[second]
    )
    return t_stat, p_value, best, second
\end{lstlisting}

The paired \textit{t}-test reports both the \textit{t}-statistic and \textit{p}-value.
If $p < 0.05$, the best model's superiority is statistically significant. If
$p \geq 0.05$, the models are statistically comparable and the choice is indifferent.
Results are printed in the notebook and logged to the experiment CSV per run.

% ═══════════════════════════════════════════════════════════════════════════════
\chapter{Advanced Analysis}
% ═══════════════════════════════════════════════════════════════════════════════

\section{Feature Importance}

\begin{figure}[H]
  \centering
  \includegraphics[width=\textwidth]{../project2_readme/decodelabs-ai-project2/assets/feature_importance.png}
  \caption{Feature importance across all models. Decision Tree uses built-in
           \code{feature\_importances\_}. KNN, SVM, and LR use permutation importance
           (10 repeats). Petal length and petal width dominate in every model —
           confirming the EDA finding. Sepal width contributes the least.}
  \label{fig:importance}
\end{figure}

\section{SHAP Explainability}

SHAP (SHapley Additive exPlanations) provides model-agnostic feature attribution at the
individual prediction level. KernelExplainer is used for SVM, as SVM has no native
\code{feature\_importances\_} attribute.

\begin{figure}[H]
  \centering
  \includegraphics[width=0.90\textwidth]{../project2_readme/decodelabs-ai-project2/assets/shap_summary_svm.png}
  \caption{SHAP summary plot for SVM (best model). Each point represents one training
           sample. High petal length values (red) strongly push predictions toward
           Virginica, while low values (blue) push toward Setosa. Sepal width has
           near-zero SHAP values, confirming its negligible contribution.}
  \label{fig:shap}
\end{figure}

\section{Calibration Analysis}

A well-calibrated model's predicted probabilities should reflect empirical accuracy —
a model that predicts 0.8 confidence should be correct 80\% of the time.

\begin{figure}[H]
  \centering
  \includegraphics[width=0.90\textwidth]{../project2_readme/decodelabs-ai-project2/assets/calibration_curves.png}
  \caption{Calibration reliability diagrams for all models. The diagonal dashed line
           represents perfect calibration. SVM (calibration deviation 0.099) shows
           slight over-confidence at mid-range probabilities. Logistic Regression
           typically achieves the best-calibrated probabilities of all classifiers.}
  \label{fig:calibration}
\end{figure}

\section{Learning Curve Analysis}

\begin{figure}[H]
  \centering
  \includegraphics[width=0.90\textwidth]{../project2_readme/decodelabs-ai-project2/assets/learning_curve_svm.png}
  \caption{Learning curve for SVM (best model). Training score (blue) and validation
           score (orange) converge as training size increases, with shaded $\pm$1 std
           deviation bands. The narrow gap between training and validation scores at
           full training size indicates a well-fitted model — neither overfitting nor
           underfitting. Adding more data would yield diminishing returns.}
  \label{fig:learning}
\end{figure}

\section{Class Imbalance Demonstration}

\sidenote{The ``accuracy mirage'' — 99\% accuracy on imbalanced data can mean the model predicts only the majority class.}

To illustrate the accuracy mirage described in the DecodeLabs brief, Versicolor is
artificially downsampled to 10 samples, creating a 50:10:50 imbalanced dataset. A KNN
classifier trained on this imbalanced set achieves high accuracy (reflecting the
majority-class dominance) but shows a dramatic collapse in F1 macro score — demonstrating
that accuracy alone is an unreliable metric for imbalanced classification problems.

This directly motivates the use of F1 macro as the primary ranking metric throughout
this project, as it treats all classes equally regardless of their support size.

% ═══════════════════════════════════════════════════════════════════════════════
\chapter{Production Deployment}
% ═══════════════════════════════════════════════════════════════════════════════

\sidenote{A sklearn Pipeline is the production unit: it chains preprocessing and inference into one serialisable object.}

\section{sklearn Pipeline Object}

The best model (SVM) is wrapped in a \code{sklearn.pipeline.Pipeline} that chains
\code{StandardScaler} and the classifier into a single object. This eliminates manual
scaling at inference time and prevents data leakage during cross-validation:

\begin{lstlisting}[caption={Production sklearn Pipeline},label={lst:pipeline}]
from sklearn.pipeline import Pipeline
from sklearn.preprocessing import StandardScaler
from sklearn.svm import SVC

pipeline = Pipeline([
    ('scaler', StandardScaler()),
    ('classifier', SVC(
        kernel='rbf',
        probability=True,
        random_state=42
    ))
])

# Fit on raw (unscaled) training data
pipeline.fit(X_train, y_train)

# Predict on raw (unscaled) test data — scaler applied internally
predictions = pipeline.predict(X_test)
\end{lstlisting}

\section{Model Persistence}

The fitted Pipeline is serialised using \code{joblib} to the \code{models/} directory.
On reload, the pipeline produces identical predictions, verified programmatically:

\begin{lstlisting}[caption={Model persistence and reload verification},label={lst:persist}]
import joblib, os

# Save
model_path = os.path.join(MODEL_DIR, 'best_model_pipeline.pkl')
joblib.dump(pipeline, model_path)

# Reload and verify
reloaded = joblib.load(model_path)
assert np.array_equal(
    pipeline.predict(X_test),
    reloaded.predict(X_test)
), "Predictions do not match — persistence failed"
print("Model persistence: PASS")
\end{lstlisting}

\section{Interactive Prediction Interface}

An \code{ipywidgets} interface provides real-time prediction directly within the Jupyter
notebook. Four float sliders control the input measurements; predicted species and class
probabilities update instantly as sliders are adjusted.

A configurable confidence threshold (default: 70\%) detects uncertain predictions. If the
maximum class probability falls below this threshold, the interface returns an
``Uncertain'' label rather than a potentially unreliable class name — a critical pattern
in production ML systems.

\begin{lstlisting}[caption={Confidence-threshold prediction function},label={lst:predict}]
def end_to_end_predict(sl, sw, pl, pw,
                       threshold=CONFIDENCE_THRESHOLD):
    """Predict with confidence threshold guard."""
    X_input = np.array([[sl, sw, pl, pw]])
    proba    = pipeline.predict_proba(X_input)[0]
    pred_idx = np.argmax(proba)
    confidence = proba[pred_idx]

    if confidence < threshold:
        return {
            'class':      'Uncertain',
            'confidence': confidence,
            'message':    f'Below {threshold*100:.0f}% threshold'
        }
    return {
        'class':         TARGET_NAMES[pred_idx],
        'confidence':    confidence,
        'probabilities': dict(zip(TARGET_NAMES, proba))
    }
\end{lstlisting}

\section{Experiment Logging}

Every notebook run appends a structured entry to \code{logs/experiment\_log.csv},
capturing all parameters and metrics for full experiment reproducibility:

\begin{table}[H]
\centering
\caption{Experiment log entry (run: 2026-05-27T20:39:06)}
\begin{tabular}{@{}L{4cm}L{9cm}@{}}
\toprule
\rowcolor{tableheader}
\textcolor{white}{\textbf{Field}} &
\textcolor{white}{\textbf{Value}} \\
\midrule
Author          & Ali Ahmad \\
\rowcolor{tablerow}
Data source     & HuggingFace \\
Data hash       & \code{67037d1d5dd80fe8...} \\
\rowcolor{tablerow}
Best model      & SVM \\
Optimal $k$     & 3 \\
\rowcolor{tablerow}
Test accuracy   & 0.9667 \\
Test F1         & 0.9666 \\
\rowcolor{tablerow}
Test precision  & 0.9697 \\
Test recall     & 0.9667 \\
\rowcolor{tablerow}
ROC-AUC         & 0.9967 \\
Weighted F1     & 0.9599 \\
\rowcolor{tablerow}
Calibration dev & 0.099 \\
\bottomrule
\end{tabular}
\end{table}

% ═══════════════════════════════════════════════════════════════════════════════
\chapter{Modular Package and Unit Testing}
% ═══════════════════════════════════════════════════════════════════════════════

\sidenote{22 unit tests cover every critical pipeline stage. A test suite is the engineering contract with your future self.}

\section{Package Architecture}

The \code{pipeline/} directory implements a clean Python package with one module per
pipeline concern:

\begin{table}[H]
\centering
\caption{pipeline/ module responsibilities}
\begin{tabular}{@{}L{4cm}L{9cm}@{}}
\toprule
\rowcolor{tableheader}
\textcolor{white}{\textbf{Module}} &
\textcolor{white}{\textbf{Responsibility}} \\
\midrule
\code{data\_loader.py}  & HF/sklearn loading, column normalisation, SHA-256 hash \\
\rowcolor{tablerow}
\code{preprocessor.py}  & X/y split, stratified split, StandardScaler, IQR outliers \\
\code{trainer.py}       & 4 base models, VotingClassifier, sklearn Pipeline \\
\rowcolor{tablerow}
\code{evaluator.py}     & 5 metrics, CV, t-test, custom scorer, calibration, CSV log \\
\code{visualizer.py}    & All 22 plot functions, returning figure objects \\
\bottomrule
\end{tabular}
\end{table}

\section{Unit Tests}

22 pytest unit tests cover all critical pipeline stages:

\begin{table}[H]
\centering
\caption{Unit test coverage summary}
\begin{tabular}{@{}L{8cm}C{2cm}@{}}
\toprule
\rowcolor{tableheader}
\textcolor{white}{\textbf{Test}} &
\textcolor{white}{\textbf{Status}} \\
\midrule
HF loader returns DataFrame with correct shape          & \cmark \\
\rowcolor{tablerow}
Column names match \code{FEATURE\_NAMES}                & \cmark \\
Target has exactly 3 unique values                       & \cmark \\
\rowcolor{tablerow}
SHA-256 hash is 64-character hex string                  & \cmark \\
Feature/target split produces correct shapes             & \cmark \\
\rowcolor{tablerow}
Stratification preserves class balance within 5\%        & \cmark \\
Train/test split is approximately 80/20                  & \cmark \\
\rowcolor{tablerow}
Scaler not fitted on test data (shape check)             & \cmark \\
All 4 base models train without error                    & \cmark \\
\rowcolor{tablerow}
All models produce predictions of correct shape          & \cmark \\
\code{compute\_metrics} returns all 5 required keys      & \cmark \\
\rowcolor{tablerow}
CV scores array length equals \code{N\_SPLITS}           & \cmark \\
\textit{t}-test returns \textit{p}-value between 0 and 1 & \cmark \\
\rowcolor{tablerow}
VotingClassifier trains and predicts correctly           & \cmark \\
sklearn Pipeline produces identical predictions          & \cmark \\
\rowcolor{tablerow}
Pipeline pkl reloads and matches original predictions    & \cmark \\
\code{detect\_outliers\_iqr} returns correct columns     & \cmark \\
\rowcolor{tablerow}
\code{validate\_shapes} raises on shape mismatch         & \cmark \\
Weighted F1 scorer is callable                           & \cmark \\
\rowcolor{tablerow}
Calibration score returns float in range [0, 1]          & \cmark \\
\code{full\_report} returns dict with correct keys       & \cmark \\
\rowcolor{tablerow}
Cross-validation returns mean and fold scores            & \cmark \\
\bottomrule
\multicolumn{2}{r}{\textbf{22 / 22 tests passing}}
\end{tabular}
\end{table}

% ═══════════════════════════════════════════════════════════════════════════════
\chapter{Problems Faced}
% ═══════════════════════════════════════════════════════════════════════════════

\section{HuggingFace Column Schema Inconsistency}

The \code{scikit-learn/iris} dataset on HuggingFace has undergone schema changes across
library versions, using at least three different column naming conventions. Building a
normalisation function robust enough to handle all three variants
(\code{SepalLengthCm}, \code{sepal.length}, \code{sepal length (cm)}) required careful
inspection of the DataFrame on first load and a multi-branch mapping strategy.

\textbf{Resolution:} The \code{\_normalise\_columns()} function checks for three known
schemas in order and falls back to positional assignment for unrecognised formats.

\section{NumPy Version Conflicts}

The HuggingFace \code{datasets} library and SHAP have conflicting NumPy version
requirements. Installing them in arbitrary order caused \code{ValueError} and
\code{ImportError} at runtime. The issue manifested as silent failures in SHAP's
KernelExplainer when NumPy was downgraded post-install.

\textbf{Resolution:} NumPy was pinned last in the installation sequence, and
\code{requirements.txt} specifies \code{numpy>=1.26.0,<2.0.0} to constrain the
acceptable range.

\section{SVM Probability Calibration Cost}

Setting \code{probability=True} on SVC triggers Platt scaling via internal
cross-validation during \code{.fit()}, substantially increasing training time. For the
Iris dataset this is negligible, but it introduced unexpected delays during development
when testing with larger synthetic datasets.

\textbf{Resolution:} Documented as a known limitation. For production with large
datasets, \code{CalibratedClassifierCV} with \code{cv='prefit'} is the appropriate
alternative.

\section{SHAP KernelExplainer Memory and Speed}

KernelExplainer is computationally expensive for SVM because it is model-agnostic
and must evaluate the model exponentially many times per prediction. Using the full
training set as background caused notebook timeouts.

\textbf{Resolution:} Background reduced to 20 randomly sampled training points
(\code{nsamples=50} perturbations per explanation). This is an approximation, documented
in the model card.

\section{ipywidgets Rendering in Exported Notebooks}

\code{ipywidgets} renders interactively inside a running Jupyter session but does not
display in static HTML or PDF exports of the notebook. Reviewers viewing the notebook on
GitHub see the cell output as empty.

\textbf{Resolution:} A static demo cell with 5 hardcoded sample predictions was added
before the widget cell, ensuring static reviewers see representative output even without
running the notebook.

% ═══════════════════════════════════════════════════════════════════════════════
\chapter{Limitations}
% ═══════════════════════════════════════════════════════════════════════════════

\begin{warningbox}[Known Limitations]
The following limitations apply to this project. They are documented here for
intellectual honesty and to frame the scope of valid conclusions.
\end{warningbox}

\section{Dataset Scale}

The Iris dataset contains 150 samples — a quantity so small that a single misclassification
changes accuracy by 0.67 percentage points. Performance differences between models are
not statistically robust at this scale. The near-identical performance of SVM, KNN, LR,
and Ensemble (all 96.67\%) should be interpreted as ``all are effective'' rather than
evidence of superiority for any particular algorithm.

\section{KNN Inference Complexity}

KNN does not learn a parametric model. At inference time, it computes the distance
between the query point and all $N$ training points — $O(N \cdot d)$ complexity. This is
acceptable for 120 training samples but becomes infeasible at industrial scale without
approximate nearest-neighbour indices (FAISS, Annoy, ScaNN).

\section{Benchmark vs Real-World Gap}

The Iris dataset is artificially clean: perfectly balanced classes, no missing values, no
outliers of consequence, and at least one class (Setosa) that is trivially linearly
separable. Real classification tasks exhibit noise, imbalance, and overlapping
distributions far more extreme than those observed here. Performance metrics achieved on
Iris do not generalise as predictions of performance on real datasets.

\section{GridSearchCV Coverage}

The hyperparameter search uses a fixed, manually defined grid. This approach misses
potentially better configurations outside the grid and is computationally wasteful for
small gains. Bayesian optimisation (Optuna, Ray Tune) would provide more efficient
coverage of a larger hyperparameter space.

\section{SHAP Approximation}

The KernelExplainer SHAP values are computed with a background of 20 samples and
\code{nsamples=50} perturbations per data point. This is a computational approximation
of the exact Shapley values. The approximation error is small for well-separated classes
but may be non-trivial near decision boundaries.

% ═══════════════════════════════════════════════════════════════════════════════
\chapter{Conclusion}
% ═══════════════════════════════════════════════════════════════════════════════

This project has successfully implemented a complete, production-grade supervised
learning pipeline for multi-class classification as part of the DecodeLabs AI
Industrial Training Program. All four formal specification requirements have been
satisfied, and the implementation substantially exceeds the minimum specification in
depth, engineering rigour, and portfolio quality.

\section{Key Findings}

\begin{enumerate}[itemsep=6pt]

  \item \textbf{SVM with RBF kernel achieves the highest performance:}
        Accuracy 96.67\%, F1 0.9666, ROC-AUC 0.9967 on the test set, and the highest
        cross-validation mean F1 with the lowest standard deviation — indicating both
        superior performance and lower generalisation variance.

  \item \textbf{Petal dimensions are the dominant features:}
        Both the Decision Tree structure and SHAP analysis confirm that petal length and
        petal width are the primary discriminating features. Sepal width contributes
        negligibly to classification accuracy.

  \item \textbf{Linear separability is class-dependent:}
        PCA and t-SNE projections confirm that Setosa is perfectly linearly separable
        from the other two species. The Versicolor--Virginica boundary is non-linear,
        which explains why the RBF kernel outperforms linear models when evaluated on
        cross-validation variance.

  \item \textbf{Accuracy alone is an insufficient metric:}
        The class imbalance demonstration showed that a model predicting only the
        majority class can achieve high accuracy while F1 collapses — validating the
        use of macro F1 as the primary evaluation metric throughout this project.

  \item \textbf{The sklearn Pipeline is the correct production unit:}
        Encapsulating preprocessing and inference in a single serialisable object
        eliminates the most common source of production ML bugs: inconsistent feature
        scaling between training and inference.

\end{enumerate}

\section{Lessons Learned}

The most valuable engineering lesson from this project is that \textit{correct process
matters more than model choice on clean data}. On the Iris benchmark, the choice of
algorithm is almost irrelevant — what distinguishes a professional implementation is
the absence of data leakage, the use of stratified splits, the proper estimation of
generalisation error through cross-validation, and the discipline to persist models in
a reproducible format with complete experiment logging.

\section{Bridge to Project 3}

Project 2 has established the supervised learning pipeline that Project 3 will extend
into deep learning and computer vision. The conceptual bridge is direct: Project 1 used
exact dictionary lookup (discrete matching); Project 2 uses learned statistical decision
boundaries on tabular features; Project 3 will use convolutional neural networks to learn
hierarchical feature representations directly from pixel data — extending the same
IPO framework to a fundamentally higher-dimensional input space.

% ═══════════════════════════════════════════════════════════════════════════════
% REFERENCES
% ═══════════════════════════════════════════════════════════════════════════════
\chapter*{References}
\addcontentsline{toc}{chapter}{References}

\begin{enumerate}[leftmargin=*, label={[\arabic*]}, itemsep=6pt]

  \item \label{fisher1936}
        Fisher, R. A. (1936). The use of multiple measurements in taxonomic problems.
        \textit{Annals of Eugenics}, 7(2), 179--188.

  \item \label{cover1967}
        Cover, T., \& Hart, P. (1967). Nearest neighbor pattern classification.
        \textit{IEEE Transactions on Information Theory}, 13(1), 21--27.

  \item \label{quinlan1986}
        Quinlan, J. R. (1986). Induction of decision trees.
        \textit{Machine Learning}, 1(1), 81--106.

  \item \label{vapnik1995}
        Vapnik, V. N. (1995). \textit{The Nature of Statistical Learning Theory}.
        Springer, New York.

  \item \label{cox1958}
        Cox, D. R. (1958). The regression analysis of binary sequences.
        \textit{Journal of the Royal Statistical Society: Series B}, 20(2), 215--242.

  \item \label{dietterich2000}
        Dietterich, T. G. (2000). Ensemble methods in machine learning.
        \textit{International Workshop on Multiple Classifier Systems}, 1--15.

  \item \label{lundberg2017}
        Lundberg, S. M., \& Lee, S.-I. (2017). A unified approach to interpreting
        model predictions. \textit{Advances in Neural Information Processing Systems},
        30, 4765--4774.

  \item \label{maaten2008}
        van der Maaten, L., \& Hinton, G. (2008). Visualizing data using t-SNE.
        \textit{Journal of Machine Learning Research}, 9, 2579--2605.

  \item \label{pedregosa2011}
        Pedregosa, F., et al. (2011). Scikit-learn: Machine learning in Python.
        \textit{Journal of Machine Learning Research}, 12, 2825--2830.

  \item \label{wolf2020}
        Wolf, T., et al. (2020). HuggingFace's Transformers: State-of-the-art natural
        language processing. \textit{Proceedings of EMNLP 2020 (System Demonstrations)},
        38--45.

  \item \label{decodelabs2026}
        DecodeLabs. (2026). \textit{Industrial Training Kit — Batch 2026: Artificial
        Intelligence Project 2 — Data Classification Using AI}.
        DecodeLabs, Greater Lucknow, India. \url{https://www.decodelabs.tech}

\end{enumerate}

\end{document}
